\quad O Shell Sort é um algoritmo de ordenação baseado no Insertion Sort,
porém utiliza comparações entre elementos que estão separados por
determinados intervalos, chamados de \textit{gaps}. Essa estratégia
permite reduzir gradualmente a desordem do vetor antes da etapa final de
ordenação.

\quad Inicialmente, é definido um valor de \textit{gap} que determina a
distância entre os elementos que serão comparados. O vetor é então
percorrido considerando essa distância, realizando uma ordenação por
inserção entre os elementos pertencentes aos diferentes grupos formados.
Após cada etapa, o valor do \textit{gap} é reduzido utilizando um divisor,
fazendo com que os elementos sejam comparados a distâncias cada vez
menores.

\quad O processo continua até que o \textit{gap} seja igual a 1. Nesse
momento, é realizada uma última ordenação semelhante ao Insertion Sort,
porém o vetor já se encontra parcialmente ordenado devido às etapas
anteriores, reduzindo bastante a quantidade de deslocamentos necessários.

\begin{lstlisting}[caption={Implementação do Shell Sort}]
void shellSort(int* vetor, int tamanho, int divisor) {

    for (int gap = tamanho / divisor; gap > 0; gap /= divisor) {
        for (int i = gap; i < tamanho; i++) { 

            int temp = vetor[i]; 

            int j = i;

            while (j >= gap && vetor[j - gap] > temp) {

                vetor[j] = vetor[j - gap]; 
                j -= gap; 
            }

            vetor[j] = temp; 
        }
    }
}
\end{lstlisting}

\quad Na primeira parte do algoritmo, o valor inicial do \textit{gap} é
obtido pela divisão do tamanho do vetor pelo divisor informado. A cada
nova etapa, o \textit{gap} é novamente dividido pelo mesmo divisor. Dessa
forma, o intervalo entre os elementos comparados diminui
progressivamente.

\quad O segundo laço percorre o vetor a partir da posição correspondente
ao \textit{gap}. Para cada elemento, a variável \texttt{temp} armazena
seu valor, enquanto o terceiro laço desloca os elementos maiores que
\texttt{temp} para a posição seguinte dentro do intervalo definido pelo
\textit{gap}.

\quad Quando não existem mais elementos maiores que \texttt{temp} na
sequência considerada, o valor é inserido em sua posição correta. Quando
o \textit{gap} chega a 1, o algoritmo realiza uma ordenação por inserção
considerando elementos consecutivos. Como as etapas anteriores já
reduziram parte da desordem do vetor, essa última etapa tende a exigir
menos deslocamentos.

\quad Inicialmente, pretendia-se permitir a alteração do divisor utilizado
na redução dos \textit{gaps} durante a execução do algoritmo, com o objetivo
de comparar o desempenho do Shell Sort utilizando diferentes sequências de
\textit{gaps}. Entretanto, devido a dificuldades na implementação dessa
alteração, foi adotado o divisor padrão igual a 2 para os experimentos
realizados neste trabalho. Dessa forma, a sequência de \textit{gaps} é
obtida pela divisão sucessiva do tamanho do vetor por 2 até que o
\textit{gap} seja igual a 1.

